% =====================================================================
%  results_c8.tex
%  Campaign 8: the candidate campaign. Opening of the section, written
%  while the campaign runs (launched 1 September 2026 on wks720, 60
%  evaluations, about 17.5 h). Paste after results_c7.tex.
%
%  Every result slot is an \openpoint. Fill them from the checkpoint
%  with inspect_front.py and c7_figures_v2.py adapted to out_c8.
%
%  Labels defined: sec:res-c8, sec:res-c8-intent, sec:res-c8-hyp,
%  sec:res-c8-doe, sec:res-c8-front, sec:res-c8-split, sec:res-c8-3d,
%  sec:res-c8-candidates, tab:c8-hypotheses, tab:c8-spec.
%  Labels used: sec:meth-c8-space, sec:meth-ctrl, sec:meth-lax,
%  sec:res-c7-screens, sec:res-c7-rodded, sec:res-c7-limits,
%  tab:ctrl-bounds, eq:ktarget-3d.
% =====================================================================

\section{Campaign 8: the candidate campaign}
\label{sec:res-c8}

Campaign 8 is the campaign whose front is proposed as the set of
candidate cores for the LABGENE-class reactor. Every earlier campaign
established something that this one fixes, and every question that
this one asks is one the earlier campaigns could not answer.
Section~\ref{sec:meth-c8-space} states the design space and
Table~\ref{tab:c8-spec} summarises the specification as launched.

\subsection{Intent}
\label{sec:res-c8-intent}

\begin{table}[htbp]
  \centering
  \small
  \caption{Campaign 8 as launched on 1 September 2026. The reactivity
  ceiling is the three-dimensional derived value of
  Table~\ref{tab:ctrl-bounds}.}
  \label{tab:c8-spec}
  \begin{tabularx}{\textwidth}{@{}>{\raggedright\arraybackslash}p{4.2cm}>{\raggedright\arraybackslash}X@{}}
    \toprule
    Item & Setting \\
    \midrule
    Design variables & $e$ (3.0 to \SI{13.913}{wt\%}), $w_\mathrm{Gd}$ (0 to \SI{8}{wt\%}), $t_\mathrm{refl}$ (2.0 to \SI{5.66}{cm}), $n_\mathrm{Gd}$ (12 to 40, ladder) \\
    Fixed & pitch \SI{1.26}{cm}, active height \SI{120}{cm}, clearance \SI{7.08}{cm} (barrel \num{5.08} plus downcomer \num{2.0}), \SI{1000}{ppm} soluble boron \\
    Objectives & maximise EFPD, minimise $F_{\Delta H}$ with all rods out \\
    Reactivity window & $1.02 \le k_\mathrm{eff}^\mathrm{core} \le 1.166$, core basis \\
    Control screen & $k_\mathrm{ARI} \le 0.99$ with RE1 to RE4 inserted, B$_4$C, constrained \\
    Second reading & $k_\mathrm{RE12}$ with RE1 and RE2 inserted, recorded only \\
    End of cycle & $k_\mathrm{target}^\mathrm{3D}$ of Equation~\eqref{eq:ktarget-3d}, 1.0839 to 1.0820 \\
    Acquisition & minimum separation 0.14, feasibility margin 1.5 standard deviations \\
    Budget & design of experiments 36, four infill blocks of six, 60 evaluations \\
    Fidelity & depletion $4000\times60$, core solves $100000\times170$ with 60 inactive \\
    \bottomrule
  \end{tabularx}
\end{table}

Three intentions distinguish this campaign from its predecessors.

\begin{enumerate}
  \item The front is meant to be read as a candidate list, so every
        design on it must be geometrically consistent, controllable by
        measurement, and evaluated against a finite-height end of
        cycle. Campaign 7 satisfied none of the three in full.

  \item The campaign measures, for every design, whether the first two
        regulating banks alone suffice. If a design is controllable
        with RE1 and RE2, the all-banks configuration that
        Section~\ref{sec:res-c7-rodded} showed to exceed the peaking
        screen is never required in operation, since the shutdown banks
        remain available for scram.

  \item With four design variables and nine design-of-experiments
        points per variable, the campaign has the densest sampling of
        the series, and it is the first in which the reflector
        thickness acts on the radial power shape alone
        (Section~\ref{sec:meth-lax}).
\end{enumerate}

\subsection{Hypotheses stated before the results}
\label{sec:res-c8-hyp}

The following expectations were written down before the design of
experiments closed, so that the results can be read against them
rather than fitted to them. Each is derived from a measurement of
Campaign 7 or of the axial study.

\begin{table}[htbp]
  \centering
  \small
  \caption{Pre-registered expectations for Campaign 8 and their basis.}
  \label{tab:c8-hypotheses}
  \begin{tabularx}{\textwidth}{@{}>{\raggedright\arraybackslash}X>{\raggedright\arraybackslash}X@{}}
    \toprule
    Expectation & Basis \\
    \midrule
    The peaking floor of the front rises from 1.525 to about 1.60
      & Five of six Campaign 7 control-screen front designs are unbuildable at pitch 1.26, and the reflector lever is capped at \SI{5.66}{cm} \\
    Cycle lengths are 450 to 500 EFPD shorter than an uncorrected run would give
      & The axial factor raises the target by \SI{2900}{pcm} at a median slope of \SI{607}{pcm} per MWd/kgHM \\
    The first two banks are worth about half of the sixteen
      & Smoke-test ratio 0.47 to 0.53 across a factor of three in enrichment \\
    Designs controllable with two banks occupy a narrow window, $k_\mathrm{eff}^\mathrm{core}$ below about 1.074
      & Two-bank worth of about \SI{8450}{pcm} at the \SI{1000}{pcm} margin \\
    The derived ceiling of 1.166 binds rarely and after the control screen
      & Of the Campaign 7 archive, 49 of 60 designs have $g_\mathrm{ctrl}$ binding first at that ceiling \\
    Every infill block moves all four variables and returns feasible designs
      & The Campaign 6 acquisition settings are restored explicitly \\
    The feasibility fraction of the design of experiments is below the \SI{31}{\%} of Campaign 7
      & The reactivity window is narrower and the enrichment box wider \\
    No design is censored and no design approaches \SI{75}{MWd/kgHM}
      & Campaign 7 maximum end-of-cycle burnup 63.8, and the corrected target shortens cycles further \\
    \bottomrule
  \end{tabularx}
\end{table}

\subsection{Design of experiments}
\label{sec:res-c8-doe}

\openpoint{Feasibility census of the 36 design-of-experiments points
against the seven recorded quantities, the count below the reactivity
floor, and the count of designs controllable with two banks. Compare
with the eighth row of Table~\ref{tab:c8-hypotheses}.}

\openpoint{Constraint activity in the plane of Figure~\ref{fig:c7-screens}
redrawn for Campaign 8, with the 1.166 ceiling and the control screen,
and the count of designs where each binds first.}

\subsection{The infill blocks and the front}
\label{sec:res-c8-front}

\openpoint{Per-block table in the form of Table~\ref{tab:c7-blocks},
with the four-variable spread of each batch. The sixth row of
Table~\ref{tab:c8-hypotheses} is decided here.}

\openpoint{The feasible front with the three-dimensional-consistent
cycle lengths, the objective-space figure in the form of
Figure~\ref{fig:c7-pareto}, the hypervolume history, and the
convergence statement under the Campaign 6 rule.}

\openpoint{Statistical error on the cycle length of every front design
by the propagation of Section~\ref{sec:res-c7-uncertainty}, and the
measured $\sigma(F_{\Delta H})$ from the core tallies, so that both
axes of the front carry an uncertainty.}

\subsection{The two-level reading}
\label{sec:res-c8-split}

\openpoint{The front split by $g_\mathrm{ctrl,12}$. For each front
design, the sixteen-bank and two-bank worths, the two rodded peaking
values, and whether two banks hold the core at the margin. State the
ratio of the two worths against the third row of
Table~\ref{tab:c8-hypotheses}.}

\openpoint{Rodded peaking in the two-bank configuration against the
screening bound of 2.0, as the operating counterpart of
Figure~\ref{fig:c7-rodded}.}

\subsection{Three-dimensional confirmation of the candidates}
\label{sec:res-c8-3d}

\openpoint{For each candidate, the finite-height core with the
benchmark hardware stack and the core barrel: unrodded and rodded
eigenvalues, $L_\mathrm{ax}$ against the campaign value of
\num{1.0289}, and the three-dimensional $F_{\Delta H}$ next to the
two-dimensional objective.}

\openpoint{Downcomer sensitivity at \SI{3.7}{cm} on the candidates.}

\subsection{Candidate designs}
\label{sec:res-c8-candidates}

\openpoint{The candidate table: each front design with its four
variables, both objectives with uncertainties, the reactivity margins,
the two-level control result, the three-dimensional confirmation, and
the location of the reported prototype enrichment figures relative to
the front.}

%% AUTHOR: the physical reading of the final front and the choice of
%% the candidate or candidates to carry into the conclusions.
